\subsection*{MT-M28. Full adelic trace class on every positive-width critical strip}
Let \(K\) be a fixed number field and retain the actual summation
image and test-function conventions of MT-M24 and MT-M27. Choose
an archimedean norm splitting \(C_K=C_K^1\times\mathbf R_+^\times\)
and give \(C_K^1\) Haar volume one. For \(a\in\mathbf R\), \(s\ge0\)
and \(I=[u,v]\), \(u<v\), put
\[
 \begin{split}
 H_{I,s}&=L^2\!\left(C_K,
    \max\{|g|^{2u},|g|^{2v}\}
         (1+\log^2|g|)^{s/2}\,dg\right),\\
 Z_{a,I,s}&=H_{I,s}/
       \overline{|\cdot|^a\Theta_K}^{\,H_{I,s}},
 \qquad \Theta_K=\Phi_K(\mathcal S_0(\mathbf A_K)).
 \end{split}                                                    \tag{M28.1}
\]
Assume GRH for all unitary Hecke characters of this \(K\), as
supplied by the conditional input \(M\), and assume
\[
                         \tfrac12-a\in(u,v).                  \tag{M28.2}
\]
Every \(f\in\mathcal D(C_K)\) then induces a trace-class
convolution operator on the Hilbert quotient (M28.1). With the
character convention of MT-M22 and with each unitary extension
\(\beta_\chi\) chosen trivial on the norm splitting, its trace is
\[
 \operatorname{Tr}_{Z_{a,I,s}}C_f
   =\sum_{\chi\in\widehat{C_K^1}}\ \sum_{\Xi_\chi(\rho)=0}
          m_{\chi,\rho}\,
          \widehat f(\beta_\chi|\cdot|^{\,\rho-a}).
                                                               \tag{M28.3}
\]
The sum is absolutely convergent, uniformly on bounded test
sets, and the trace norms are bounded on those sets. At every
fixed compact-invariance and compact-support test level the
map \(f\mapsto C_f\) is continuous into trace norm.

The denominator in (M28.1) is exactly the kernel of all
character/zero jets
\[
 \partial_z^k\widehat h(\beta_\chi|\cdot|^z)
                  \big|_{z=\rho-a},\qquad 0\le k<m_{\chi,\rho}.
                                                               \tag{M28.4}
\]
Consequently the corrected full smooth quotient
\[
 |\cdot|^a\mathcal S_{\mathrm E}(C_K)\big/
           |\cdot|^a\overline{\Theta_K}^{\,\mathcal S_{\mathrm E}(C_K)}
                                                               \tag{M28.5}
\]
maps continuously and injectively into \(Z_{a,I,s}\), with
proper dense image. In particular, at \(a=1/2\), every
\(I=[-\alpha,\alpha]\) with \(\alpha>0\) is allowed.
The zero-width completion remains governed by MT-M25.

\paragraph{Classical input and scope of the uniform proof.}
The analytic input besides \(M\) is the usual Hecke
continuation and functional equation, with its finite-order
property, already used in NT-L03, NT-Z10 and NT-Z17.
The continuation/functional-equation locator is
Ramakrishnan--Valenza, \emph{Fourier Analysis on Number Fields},
Theorem~7-19, printed pp.~279--280; the real and complex
local factors are in \S7.1, printed pp.~244--245, with their
explicit integral calculations on pp.~250--252.
These local formulas and the entire supplied proof of
Theorem~7-19 were read, and the relevant local-factor pages
were inspected as original page images. We do not claim that
Theorem~7-19 separately states a uniform finite-order estimate:
only the classical finite-order property of each function is
used below, and all uniform estimates are proved here.
The source is
\url{https://webusers.imj-prg.fr/~julien.marche/Ref_GL_GT/Fourier_Analysis_on_Number_Fields.pdf}.
The exact summation-image identification is MT-M24, and
the already proved one-character Hilbert theorem is MT-M26.
Meyer's uniform nuclear lemma from MT-M27 is not needed in
the uniform estimates below.

\begin{proof}
\emph{1. Character parameters at one elementary test level.}
Fix a compact subgroup \(k\subset C_K^1\) such that
\(T=C_K^1/k\) is a compact abelian Lie group. Its dual is
\[
                  \Lambda=\widehat T\simeq\mathbf Z^d\times F,
                                                               \tag{M28.6}
\]
where \(F\) is finite. Every bounded subset of
\(\mathcal D(C_K)\) is contained in one such invariance
level, with a common compact support in \(x=\log|g|\)
and bounded derivatives of every order on \(T\times\mathbf R\).
This is the elementary-level definition of the Bruhat test
bornology/topology used in MT-M27. We prove uniform estimates
at this level, without replacing its infinitely many
archimedean characters by a finite set.

Choose generators of the free part in (M28.6), and extend
each character to \(C_K\) by making it trivial on the chosen
norm splitting. Then the extension at lattice index
\(\nu\in\mathbf Z^d\times F\) is the product of the fixed
generator extensions and one of finitely many torsion
extensions. Each fixed extension has finite conductor.
The conductor of a product divides the least common
multiple of the conductors of its factors. Hence all finite
conductors at this level divide one fixed ideal.

For the primitive Hecke function \(L_\nu\) in the convention
of MT-M22, write its completion as
\[
 \Xi_\nu(z)=Q_\nu^{z/2}
      \prod_{j=1}^{n_K}\Gamma_{\mathbf R}(z+\mu_{\nu,j})L_\nu(z),
 \qquad \Gamma_{\mathbf R}(z)=\pi^{-z/2}\Gamma(z/2),
                                                               \tag{M28.7}
\]
up to a nonzero scalar. Here \(n_K=[K:\mathbf Q]\).
For every \(\nu\ne0\) this completion is entire. The
single trivial compact character, with the chosen extension,
gives \(\zeta_K\); for it use the pole-removed completion
and the fixed-character result MT-M26 throughout. A pure
norm character trivial on the norm splitting is trivial,
so there are no other polar exceptions in this family.

The local formulas in the cited pages show that real
places contribute shifts \(\epsilon+i\tau\), with
\(\epsilon=0\) or \(1\), and complex places contribute
\[
      |m|/2+i\tau,\qquad |m|/2+1+i\tau.
 \]
Indeed
\(\Gamma_{\mathbf C}(z)=
\Gamma_{\mathbf R}(z)\Gamma_{\mathbf R}(z+1)\);
an alternative constant normalization cancels in all
ratios below. Angular integers and real norm parameters
of products grow at most linearly in the exponent vector.
Thus, writing \(H_\nu=2+|\nu|\) for any fixed lattice norm,
\[
 \begin{gathered}
  1\le Q_\nu\le Q_0,\qquad
  \operatorname{Re}\mu_{\nu,j}\in\tfrac12\mathbf Z_{\ge0},
  \qquad |\mu_{\nu,j}|\le C H_\nu .
 \end{gathered}                                                 \tag{M28.8}
\]
The constants depend only on the field and this test level.
They are independent of \(\nu\).

\emph{2. Uniform gamma ratios and finite-\(L\) bounds.}
We record the elementary gamma estimates needed to avoid
an uncanceled, possibly very large, completed-function
normalization. For \(\operatorname{Re}z\ge c>0\), the
digamma series gives
\[
 \psi(z)=-\gamma+
      \sum_{m\ge0}\left(\frac1{m+1}-\frac1{m+z}\right),
 \qquad
 |\psi(z)|\le C_c+2\log(2+|z|),
 \qquad |\operatorname{Im}\psi(z)|\le C_c.                    \tag{M28.9}
\]
Here is a direct bound. Split the series at
\(N=\lceil2|z|+2\rceil\). The finite part is bounded by
two harmonic sums, using \(|m+z|\ge m+c\); the tail is
bounded by
\(\sum_{m\ge N}|z-1|/((m+1)|m+z|)=O(1)\).
For \(z=x+iy\), the imaginary part is bounded by
\[
 |y|\sum_{m\ge0}\frac1{(m+x)^2+y^2}
 \le \frac1{2c}+\int_0^\infty
           \frac{|y|}{(t+x)^2+y^2}\,dt
 \le \frac1{2c}+\frac\pi2,
 \]
with the \(y=0\) case immediate.
Consequently a bounded change in real part, within a
right half-plane, changes \(|\Gamma_{\mathbf R}|\) by
at most a fixed power of \(2+|z|\), in either direction.
A change of imaginary part by \(h\) changes it by at
most \(e^{C_c|h|}\). These statements follow by
integrating the real or imaginary derivative of
\(\log|\Gamma_{\mathbf R}|\); the half-plane is free
of gamma poles and zeros.

Fix \(A>0\). The Euler product gives
\[
 \zeta_K(2)^{-1}\le |L_\nu(2+it)|\le\zeta_K(2),
 \qquad |L_\nu(A+2+it)|\le\zeta_K(2).
                                                               \tag{M28.10}
\]
The lower bound follows factor by factor from
\((1+x)^{-1}\ge1-x\), \(0\le x<1\), including the
omitted ramified factors. On \(\sigma=-A\), the
functional equation expresses \(L_\nu(\sigma+it)\)
as \(L_{\bar\nu}(1-\sigma-it)\), a bounded conductor
power and a product of ratios
\[
  \frac{\Gamma_{\mathbf R}(1-\sigma+
                   \operatorname{Re}\mu_{\nu,j}
                   +i(t+\operatorname{Im}\mu_{\nu,j}))}
       {\Gamma_{\mathbf R}(\sigma+\mu_{\nu,j}+it)}
                                                               \tag{M28.11}
\]
in absolute value. Conjugation was used in the
numerator; the root number has modulus one.
Use
\(\Gamma_{\mathbf R}(z+2)=
 (z/(2\pi))\Gamma_{\mathbf R}(z)\)
to move the denominator to a fixed right half-plane.
The reciprocal of the original denominator becomes
the reciprocal of that shifted gamma function times
a polynomial of fixed degree. This remains valid
at poles of the original denominator, by analytic
continuation. Now (M28.9) and (M28.8) bound every
ratio by a fixed power of \(H_\nu+|t|\).
It follows that both boundary lines of the strip
\([-A,A+2]\) satisfy such a polynomial bound.

For explicitness, the Phragm\'en--Lindel\"of step is
uniform despite using only finite order separately
for each \(\nu\). Divide \(L_\nu(z)\) by
\((H_\nu+C_A+z)^B\), with integer \(B\) sufficiently
large and \(C_A>A+2\). This denominator has no zero
in the strip, and on its boundaries its modulus
is comparable with \((H_\nu+|t|)^B\). The quotient
is uniformly bounded on both boundaries and has
finite order for each \(\nu\). Multiply it by
\[
          \exp\{-\epsilon\cos(\lambda(z-z_0))\},
 \]
where \(z_0\) is the real midpoint of the strip and
\(\lambda>0\) is small enough that
\(\cos(\lambda(\sigma-z_0))>c>0\) throughout it.
The multiplier has modulus at most one on the
vertical boundaries and tends to zero faster than
every finite-order growth on horizontal boundaries
as their heights tend to infinity. The maximum
principle on rectangles, followed first by their
height tending to infinity and then by
\(\epsilon\downarrow0\), gives the same uniform
bound throughout the strip. Thus, for every
fixed \(A\),
\[
       |L_\nu(\sigma+it)|
           \le C_A(H_\nu+|t|)^{B_A}
                 \quad (|\sigma|\le A,\ \nu\ne0).
                                                               \tag{M28.12}
\]
No GRH has been used in (M28.8)--(M28.12).

Jensen's formula, centered at \(2+iT\), now gives
for each fixed radius \(R\)
\[
  \#\{L_\nu\text{ zeros in }|z-(2+iT)|\le R\}
             \le C_R\log(H_\nu+|T|),                           \tag{M28.13}
\]
counting multiplicity. Indeed use a larger concentric
circle of radius at least \(2R\), (M28.12) for its
upper bound and the lower bound (M28.10) at the
center. A boundary avoiding zeros can be chosen in
a fixed bounded radius interval. The logarithmic
ratio in Jensen's inequality stays bounded below.
This counts trivial zeros as well as nontrivial
ones and is uniform at every height, not only
above a character-dependent exceptional range.

\emph{3. Uniform narrow block contours.}
Continue with \(\nu\ne0\); the single polar family
member is included using MT-M26 at the end of step~4.
Choose \(0<\delta<1/4\) so small that
\[
 [b_-,b_+]=[1/2-\delta,1/2+\delta]
                         \subset a+(u,v).
 \]
By the assumed GRH, all zeros of the entire
completion \(\Xi_\nu\) lie on \(\operatorname{Re}z=1/2\).
All other zeros of the finite \(L_\nu\) are the
trivial zeros canceling gamma poles, and have
nonpositive real part. Thus all nontrivial zeros
are separated by \(\delta\) from the two vertical
edges, and all trivial zeros are separated by at
least \(b_->1/4\) from every point of the narrow
closed strip.

Put \(X_{\nu,n}=H_\nu+|n|\). Count the nontrivial
zero ordinates in \([n-2,n+3]\), with multiplicity,
by \(M_{\nu,n}\). Equation (M28.13), with a fixed
sufficient radius, gives \(M_{\nu,n}\le C\log X_{\nu,n}\).
The elementary averaging argument of MT-M18 now
chooses
\[
        h_{\nu,n}\in(n+1/4,n+3/4)
 \]
avoiding every zero ordinate, such that
\[
 \prod_{\Xi_\nu(\rho)=0}
      \min(1,|h_{\nu,n}-\operatorname{Im}\rho|)^{m_\rho}
           \ge e^{-5M_{\nu,n}}\ge X_{\nu,n}^{-C}.             \tag{M28.14}
\]
For completeness, each function
\(\log^+(1/|h-\gamma|)\) has integral at most \(2\)
on an interval of length \(1/2\). Their sum
therefore has a value at most \(5M_{\nu,n}\);
the slack gives a positive-measure set when
\(M_{\nu,n}>0\). Every factor differing from one
comes from the specified ordinate window.

Let \(\Omega_{\nu,n}\) be the rectangle with
vertical edges \(b_\pm\) and horizontal edges
\(h_{\nu,n},h_{\nu,n+1}\); write
\(\Gamma_{\nu,n}=\partial\Omega_{\nu,n}\).
Its perimeter is uniformly bounded. The disk
Blaschke/Harnack proof of MT-M18, now using the
uniform bounds (M28.10)--(M28.13), gives
\[
        \inf_{w\in\Gamma_{\nu,n}}|L_\nu(w)|
                         \ge c X_{\nu,n}^{-B}.               \tag{M28.15}
\]
Here are the details needed for uniformity. Center
a disk at \(2+i\operatorname{Im}w\), with a
zero-free boundary of radius in a fixed interval
large enough that the narrow strip point lies in
the inner half disk. Remove all finite-\(L\)
zeros in this disk by its finite Blaschke product.
The remaining zero-free factor has modulus at
most \(C X_{\nu,n}^B\) on the boundary and at least
\(\zeta_K(2)^{-1}\) at the center. Harnack applied
to the logarithm of this upper bound divided by
its modulus gives a lower bound \(cX_{\nu,n}^{-B}\)
on the inner half disk. There are at most
\(C\log X_{\nu,n}\) Blaschke factors.

On a vertical edge every nontrivial factor has
distance at least \(\delta\) from its zero, and
every trivial factor distance at least \(b_-\).
Each normalized factor is bounded below by this
distance divided by twice the disk radius.
On a horizontal edge, use (M28.14) for the
nontrivial factors, and \(b_-\) for the trivial
factors. Omitting factors from the product in
(M28.14) only increases that product.
These lower bounds, each with \(O(\log X_{\nu,n})\)
factors, prove (M28.15).
The argument includes small heights and trivial
zeros whose ordinates grow with \(\nu\). It
does not absorb any such range into a
character-dependent compactness constant.

\emph{4. Uniform lifts with the completed gamma factors canceled.}
Let
\[
 J_{\nu,n}=A(\Omega_{\nu,n})/
       \{F:F^{(j)}(\rho)=0\
             \text{for every block zero and }j<m_\rho\},
                                                               \tag{M28.16}
\]
with the quotient of the supremum norm on the
closed rectangle. This is the same finite
Banach algebra as in MT-M18. Its dimension is
\[
                  d_{\nu,n}\le C\log X_{\nu,n}.                \tag{M28.17}
\]
If the block is empty its quotient is the zero
space. All assertions below then hold with
the corresponding maps zero.

Use the entire difference-kernel lift of MT-M18,
with its Gaussian parameter fixed to \(1\):
\[
 H_{\nu,n}F(z)=\frac1{2\pi i}
       \int_{\Gamma_{\nu,n}}
       F(w)\frac{\Xi_\nu(z)-\Xi_\nu(w)}
                    {\Xi_\nu(w)(z-w)}
                         e^{(z-w)^2}\,dw.                    \tag{M28.18}
\]
As proved there, it is entire, has the prescribed
jets in this block and zero full jets in all
other blocks, and depends only on the class of
\(F\) in \(J_{\nu,n}\). For each fixed \(\nu\),
the MT-M18 construction and the contour
adaptation in MT-M26 show that it is the
Mellin transform of an element of \(\mathcal E\).
We now prove a uniform estimate for that element.

For an integer \(N\ge0\), take the two inversion
lines
\[
                 \operatorname{Re}z=\pm(N+9/4).              \tag{M28.19}
\]
They lie outside all the rectangles. On them the
Cauchy integral of the \(-\Xi_\nu(w)\) term in
(M28.18) is zero. Only the ratio
\(\Xi_\nu(z)/\Xi_\nu(w)\) remains.
Move each numerator gamma factor
\(\Gamma_{\mathbf R}(z+\mu_{\nu,j})\) to a right
half-plane by a fixed even shift, if necessary.
Every recurrence denominator has real part in
\[
        \pm(N+9/4)+\tfrac12\mathbf Z+2\mathbf Z,
 \]
and therefore absolute value at least \(1/4\).
There is no gamma pole anywhere on either
entire inversion line, for any \(\nu\) or height.
The shifted numerator and the denominator
\(\Gamma_{\mathbf R}(w+\mu_{\nu,j})\) both lie
in a fixed right half-plane.

Use (M28.9) first to change the imaginary part,
then to change the real part by their bounded
difference. Uniformly for \(z=\pm(N+9/4)+it\)
and \(w\in\Gamma_{\nu,n}\), this gives
\[
 \left|
  \frac{\Gamma_{\mathbf R}(z+\mu_{\nu,j})}
       {\Gamma_{\mathbf R}(w+\mu_{\nu,j})}
 \right|
 \le C_N(H_\nu+|t|)^{B_N}
                      e^{C_N|t-\operatorname{Im}w|}.         \tag{M28.20}
\]
The fixed recurrence factors have been included
in \(C_N\). The finite conductor ratio is bounded
by (M28.8). The finite-\(L\) numerator is bounded
by (M28.12) and its denominator by (M28.15).
Thus (M28.7) gives the uniform completed ratio
\[
 \left|\frac{\Xi_\nu(z)}{\Xi_\nu(w)}\right|
 \le C_N X_{\nu,n}^{B_N}
            (H_\nu+|t|)^{B_N}
                      e^{C_N|t-\operatorname{Im}w|}.         \tag{M28.21}
\]
The normalization scalar in (M28.7) cancels
exactly. No bound for an uncanceled completed
function, and in particular no polynomial
bound for a potentially factorial-size gamma
normalization, is asserted or used.

The Gaussian in (M28.18) has modulus
\(\exp((\operatorname{Re}z-\operatorname{Re}w)^2
       -(t-\operatorname{Im}w)^2)\).
The real-part difference is bounded for fixed
\(N\), and
\(|\operatorname{Im}w-n|\le7/4\). The Gaussian
absorbs the exponential linear factor in
(M28.21) and every polynomial in \(|t-n|\).
Also \(H_\nu+|t|\le X_{\nu,n}(1+|t-n|)\).
The denominator \(|z-w|\) is bounded below
by a positive constant on the inversion lines.
Since the contour perimeter is bounded, we obtain
\[
 |H_{\nu,n}F(\pm(N+9/4)+it)|
 \le C_N X_{\nu,n}^{B_N}\|F\|_\infty
                              e^{-c(t-n)^2},                 \tag{M28.22}
\]
with \(c>0\) independent of \(\nu,n,N\), after
reducing it once.

Invert the Mellin transform on these lines.
For \(x\ge0\) use the positive line and for
\(x\le0\) the negative line. The \(j\)-th
derivative costs \(|z|^j\), whose Gaussian
integral is bounded by a constant times
\((2+|n|)^j\). Since the resulting function
is already in \(\mathcal E\) for each
individual \(\nu\), the contour shifts are
legitimate; (M28.22) now supplies their
uniform constants. Taking the infimum over
all representatives \(F\) of a fixed block
class yields linear lifts
\(\mathcal T_{\nu,n}:J_{\nu,n}\to\mathcal E\)
with
\[
 p_{N,j}(\mathcal T_{\nu,n}b)
       \le C_{N,j}X_{\nu,n}^{B_{N,j}}\|b\|_{J_{\nu,n}},
 \qquad R_{\nu,m}\mathcal T_{\nu,n}=\delta_{mn}\operatorname{id}.
                                                               \tag{M28.23}
\]
Here \(p_{N,j}(h)=\sup_x e^{N|x|}|h^{(j)}(x)|\).
The last identity means the identity in the
matching block and zero in the others.
The exceptional \(\nu=0\) has the corresponding
fixed-character polynomial bounds by MT-M26;
enlarging the constants once includes it.
These are stable quotient-block estimates,
with no separation condition on zeros inside
one rectangle.

\emph{5. The actual quotient and its full zero-jet kernel.}
The compact character projectors are orthogonal
contractions on \(H_{I,s}\). The denominator
of (M28.1) is invariant under them: MT-M24
identifies \(\Theta_K\) with the intersection
of the two summation spaces, and compact
averaging preserves that intersection. Applying
a projector to an approximating sequence
shows that its intersection with one character
sector is exactly the Hilbert closure of the
corresponding sector image.

Apply MT-M26 in each sector, using (M28.2).
It identifies that closure with all the
sector zero jets. Every such jet is a
continuous functional on the full Hilbert
space, by the compact projection and the
interior-strip Cauchy--Schwarz estimate.
Finite character truncations converge
strongly in the Hilbert direct sum.
Therefore the closed denominator is exactly
the simultaneous kernel (M28.4).

The full smooth closed-kernel identity of
MT-M27 was proved by compact Fourier
truncation and MT-M24, independently of
the Hilbert strip width. It continues to
apply here. It proves faithfulness and
continuity of (M28.5); compactly supported
smooth radial functions and finite
character sums prove density. If this
map were onto, applying any one compact
character projector would contradict
the proper-image assertion of MT-M26.
This proves the stated proper dense
smooth inclusion without equating
the smooth and Hilbert topologies.

\emph{6. Trace-norm summation over characters and blocks.}
Let \(\pi:H_{I,s}\to Z_{a,I,s}\) be the quotient map.
The block restriction map, after compact
character projection, is
\[
  r_{\nu,n}[h]
       =\left[w\longmapsto
          \widehat h(\beta_\nu|\cdot|^{\,w-a})\right]
           \in J_{\nu,n}.
 \]
The real parts \(w-a\) lie in a fixed compact
subinterval of \((u,v)\). Cauchy--Schwarz
against the radial weight, and the fact
that the compact character has modulus
one, therefore give
\[
                \|r_{\nu,n}\|\le C                         \tag{M28.24}
\]
uniformly in both indices. As in MT-M26,
this follows by bounding the quotient
norm by the supremum of the displayed
representative on the contour.

For \(b\in J_{\nu,n}\), define
\[
  s_{\nu,n}b
       =\pi\left[
          \beta_\nu(g)^{-1}|g|^a
               (\mathcal T_{\nu,n}b)(\log|g|)
              \right].
 \]
A single fixed seminorm \(p_{N,0}\), with
\(N>|a|+\max(|u|,|v|)+1\), bounds its
Hilbert norm. Equation (M28.23) gives
\[
 \|s_{\nu,n}\|\le C X_{\nu,n}^B,\qquad
 r_{\mu,m}s_{\nu,n}
      =\begin{cases}\operatorname{id},&(\mu,m)=(\nu,n),\\
                      0,&\text{otherwise}.
        \end{cases}                                          \tag{M28.25}
\]
No compact-direction derivatives occur
in this Hilbert norm, so no extra
character factor is concealed here.

For a test \(f\) invariant under \(k\), let
\[
 f_\nu(x)=\int_{C_K^1}
          f(c\,\tau(e^x))\,\beta_\nu(c)\,dc,\qquad
 b_{\nu,n}(f)=
       \sup_{w\in\Gamma_{\nu,n}}
                      |\mathcal M f_\nu(w-a)|.
 \]
For every \(A>0\),
\[
                   b_{\nu,n}(f)\le C_A(f)X_{\nu,n}^{-A},       \tag{M28.26}
\]
where \(C_A(f)\) is bounded by finitely
many seminorms at this test level and
is uniformly bounded on bounded test
sets. To see this directly, integration
by parts on the torus, using
\((1-\Delta_T)^r\), gives arbitrary
negative powers of \(1+|\nu|\).
Integration by parts in \(x\), on the
common compact support, gives arbitrary
negative powers of \(1+|\operatorname{Im}w|\),
uniformly for the real part in the
fixed contour interval. The derivatives
of \(e^{(\operatorname{Re}w-a)x}f\)
are bounded by finitely many test
seminorms. Since
\(|\operatorname{Im}w-n|\le7/4\),
combining these two arbitrary powers
gives (M28.26). The finite group in
(M28.6) only changes a constant.

Let \(M_{f,\nu,n}\) be multiplication
in \(J_{\nu,n}\) by
\([\mathcal Mf_\nu(w-a)]\). Its norm
is at most \(b_{\nu,n}(f)\), by the
submultiplicative quotient algebra
norm. The finite-rank operator
\[
                 A_{\nu,n}(f)
                    =s_{\nu,n}M_{f,\nu,n}r_{\nu,n}
 \]
has rank at most \(d_{\nu,n}\). Hence,
as an operator on the Hilbert quotient,
\[
 \|A_{\nu,n}(f)\|_1
       \le C\,d_{\nu,n}X_{\nu,n}^B b_{\nu,n}(f).
                                                               \tag{M28.27}
\]
This is the elementary Hilbert estimate
\(\|A\|_1\le\operatorname{rank}(A)\|A\|\);
the intervening finite-dimensional
spaces need not carry Hilbert norms.
Equations (M28.17) and (M28.26) show
that
\[
                  \sum_{\nu\in\Lambda}\sum_{n\in\mathbf Z}
                                      A_{\nu,n}(f)            \tag{M28.28}
\]
converges in trace norm. For example,
take \(A>B+d+3\) in (M28.26), since
the lattice plus the height index
has dimension \(d+1\). Its tails
converge uniformly on bounded test
sets, and the resulting trace norm
is bounded by finitely many test
seminorms.

This is the actual convolution, not
only a formal operator with the
same eigenvalues. Convolution by
\(f\) is bounded on \(H_{I,s}\):
the translation estimate (M26.10)
integrated over its compact support
proves this. It preserves the closed
denominator, since
\[
 C_f(|\cdot|^a h)
       =|\cdot|^a C_{|\cdot|^{-a}f}h
 \]
and the summation image is a
convolution submodule.
On a single-character smooth vector
the operator (M28.28) equals this
convolution by the already proved
one-character synthesis and dense
extension in MT-M26. On characters
not in \(\Lambda\), both operators
are zero: \(k\)-invariance of \(f\)
makes its Fourier coefficient zero.
Finite-character smooth vectors
are dense in the full Hilbert
quotient. Both operators are
bounded, so their equality extends
to that quotient. This completes
the full adelic trace-class proof.

\emph{7. Trace, uniform convergence, and the analytic scope.}
By (M28.25), the finite-rank trace
of \(A_{\nu,n}(f)\) is the trace
of multiplication by
\([\mathcal Mf_\nu(w-a)]\) in
\(J_{\nu,n}\). The local algebra
decomposes as the direct sum of
\(\mathbf C[z]/(z-\rho)^{m_\rho}\)
over its block zeros. Multiplication
is triangular in the power bases,
with diagonal value
\(\mathcal Mf_\nu(\rho-a)\) repeated
\(m_\rho\) times. Thus this trace
is the corresponding finite part
of (M28.3). Trace-norm convergence
allows summation of these traces.
The individual spectral terms
are absolutely summable as well:
\[
 \sum_{\rho\text{ in block}}
    m_\rho|\mathcal Mf_\nu(\rho-a)|
            \le d_{\nu,n}b_{\nu,n}(f),
 \]
by the maximum principle on the
block. Equations (M28.17) and
(M28.26) give uniform tails of
this double series directly.
This proves all the asserted
trace and continuity properties.

For a specified bounded test set
at the level \(k\), its trace-class
and trace-formula conclusions use
GRH only for the characters in
\(\widehat{C_K^1/k}\). Other
character sectors are annihilated
by these tests. The global
faithfulness statement (M28.4)
uses all character sectors.
Thus the full conditional input
\(M\) is sufficient, while the
operator conclusion at one level
has the stated smaller Hecke
family as its precise analytic
input.

At \(a=1/2\), define
\(h^*(g)=\overline{h(g^{-1})}\).
The group is abelian and
unimodular. Under the assumed
GRH, every
\(\beta_\chi|\cdot|^{\rho-1/2}\)
is unitary, and (M28.3) gives
\[
       \operatorname{Tr} C_{h*h^*}
          =\sum_{\chi,\rho}
                 m_{\chi,\rho}
                 |\widehat h(
                    \beta_\chi|\cdot|^{\,\rho-1/2})|^2
          \ge0.                                               \tag{M28.29}
\]
This is a trace identity. It
does not identify the group
involution with the operator
adjoint for the specified
Hilbert inner product.
All higher zero jets remain
present in (M28.4), consistently
with the obstruction in MT-M20
to a faithful positive
full-jet realization when
multiple zeros occur.
\end{proof}

\paragraph{Run and geometric crosswalk.}
This is a new uniform transfer proof in
\texttt{GRH\_NIGHT/complexity}, following the fixed-character
MT-M18/MT-M24/MT-M26 and the full wide-strip MT-M27.
Its literal source targets are Masullo
Definition~9.2.2.31, original PDF p.~416, the Hilbert
closure in Definition~10.1.1.1, p.~434, and the
strip quotient in Definition~10.4.1.3, p.~456,
with the centering twist from MT-M24. It closes
the all-character positive-width strip gate
conditionally on \(M\); it is not a deduction
of \(M\). In the attached geometric programme,
it supplies an actual trace-class realization
for the idele/test-function quotient interface
of the \texttt{cote1}/\texttt{SA23geo-H-lidskii}
crosswalk. The positivity statement remains
the conditional trace statement (M28.29);
the stronger faithful positive-inner-product
gate is still exactly MT-M20.
